\documentclass[11pt,letterpaper]{article}
\usepackage[T1]{fontenc}
\usepackage{lmodern,amsmath,amssymb,amsthm,microtype}
\usepackage[margin=1in]{geometry}
\usepackage[hidelinks]{hyperref}
\newtheorem{theorem}{Theorem}
\title{Five-point extensions of \(G_{449}\) are five-colorable}
\author{Aafi Khan Malek\\ \small Independent researcher, Ahmedabad, India\\
\small\href{mailto:aafimalek2032@gmail.com}{aafimalek2032@gmail.com}}
\date{28 November 2026}
\begin{document}
\maketitle
Let \(P\) be the 449 points from Heule and Przybocki's
\href{https://arxiv.org/abs/2611.28518}{\(G_{449}\) construction}.
We prove an extension property for this fixed point set using
exact finite arithmetic and a stored coloring.

\begin{theorem}
For any set \(X\) of at most five points in the Euclidean plane,
the induced unit-distance graph on \(P\cup X\) is five-colorable.
\end{theorem}

Use the grid \(L=\tfrac17\mathbb Z[\exp(2\pi i/42)]\). The
\href{https://github.com/Aafimalek/novel-discoveries/blob/1c624e99b4c05f24b89d26ee941a62d7bd959fc0/hadwiger-nelson/verification/g449/prime-seven-catalogue/four-five-point-extensions/paper/four-point-extensions.pdf}
{previous calculation} covers any four added points and any five
with at least one point outside \(L\). Thus we need five distinct
new grid points. In any failure, each new point must have degree at
least five in the full extension: delete one, apply the four-point
result, and put it back in a color unused by its neighbors whenever
it has at most four neighbors.

\section{A point with one original neighbor}
The degree condition excludes every point with no original
neighbors. A new point with exactly one original neighbor
must meet all four of the other new points at unit distance.
There cannot be two such points, since they would have
three distinct common unit neighbors. Accordingly the other
four points lie on its unit circle and each has at least two
original neighbors.

Generating new points from all 546 signed grid unit directions
gives 222,264 candidates for the point of original degree one
and 7,056 candidates for its neighbors. There are 109,683
incidences between these two candidate sets. Enumerating every
four-element subset of a degree-one candidate's unit neighbors
gives 872,165 choices. Exact directions determine every edge
among the four neighbors as well. The degree condition leaves
6,839 configurations. For each, the checker finds a proper
list coloring of the five new vertices with one of 18
previously checked base colorings. Its optional second check
compares the list tests to direct Cartesian assignment enumeration.

\section{Every point has at least two original neighbors}
The 7,056 candidates have 33,908 exact grid unit edges. For each
of 37 checked base colorings, a candidate's list contains the
colors absent from its original neighbors. Each connected
component of the graph induced by the new points has a bit
mask indicating which base colorings permit a proper coloring
of that component. For a fixed partition into components the
intersection of these masks is nonempty exactly when the bank
colors the five-point extension.

The checker builds connected candidates by growing connected
sets from their least vertex and pruning only when an existing
vertex cannot attain degree five even with all available new
neighbors. A fixed parent rule removes the greatest non-root
vertex whose deletion preserves connectivity; a spanning-tree
leaf guarantees such a vertex. This generates each of the
1,119,839 qualifying connected five-sets exactly once.
Disconnected cases use their smaller connected components:
group each component size by its mask, expand only combinations
with empty mask intersection, and then check that the components
are distinct, disjoint and have no edges between them.

In component-size order
\[
  1+1+1+1+1,\quad 2+1+1+1,\quad 2+2+1,\quad
  3+1+1,\quad 3+2,\quad 4+1,\quad 5,
\]
the numbers of residual sets are respectively
91,936, 42,572, 1,398, 1,969, 36, 93, and 5.
An independent enumeration uses a deduplicated connected
frontier and direct mask joins and agrees on all 138,009
residual sets. Their union has just 348 new vertices.
Together with \(P\) this is an induced graph on 797 vertices
and 4,532 edges. One stored five-coloring covers every
residual set. A separate Laurent-polynomial calculation
checks all 317,206 pairwise distances and independently
verifies the coloring. This completes the proof.

\paragraph{Reproducibility and scope.}
The companion package pins the original graph inputs and
supplies the complete deterministic check and the color vector.
The replay uses no SAT solver and no floating-point test.
The proof has not been formalized in Lean. The result implies
that a six-chromatic induced extension of this specific
point set would need at least six new points; it does not
change the bounds for the chromatic number of the plane.

\paragraph{Contributions.}
Aafi Khan Malek selected and directed the research. ChatGPT
assisted with proof development, software and writing.
Codex was not used. This note has not been peer reviewed.
\end{document}
